\section{La configuración ganadora en los dos modelos}
\label{sec:mejor}

El último paso aplica la configuración ganadora del estudio de ablación
--- \texttt{lineal}: cabeza $\mathrm{Linear}(768 \to C)$, \textit{encoder}
entrenable --- a BERT y a DistilBERT con el mismo protocolo y los mismos
datos.

\begin{table}[h]
\centering
\small
\begin{tabular}{llccccc r}
\toprule
\textbf{Dataset} & \textbf{Modelo} & \textbf{Accuracy} & \textbf{Precision} &
\textbf{Recall} & \textbf{F1} & \textbf{F1 val.} & \textbf{Train (min)} \\
\midrule
SST-2   & BERT       & \best{0,9289} & 0,9294 & 0,9286 & \best{0,9288} & 0,9541 & 4,2 \\
SST-2   & DistilBERT & 0,9094 & 0,9099 & 0,9091 & 0,9093 & 0,9483 & \best{1,7} \\
\midrule
AG News & BERT       & \best{0,9453} & 0,9456 & 0,9453 & \best{0,9453} & 0,9481 & 9,0 \\
AG News & DistilBERT & 0,9442 & 0,9443 & 0,9442 & 0,9442 & 0,9462 & \best{5,4} \\
\midrule
Yelp    & BERT       & \best{0,9653} & 0,9653 & 0,9653 & \best{0,9653} & 0,9661 & 16,0 \\
Yelp    & DistilBERT & 0,9610 & 0,9610 & 0,9610 & 0,9610 & 0,9603 & \best{8,2} \\
\bottomrule
\end{tabular}
\caption{Configuración \texttt{lineal} en los dos modelos. BERT tiene
\num{109,5}~M de parámetros y DistilBERT \num{66,4}~M.}
\label{tab:mejor}
\end{table}

Con la cabeza lineal, BERT gana en los tres datasets: por \num{1,95} puntos en
SST-2, \num{0,43} en Yelp y \num{0,11} en AG News. Es el mismo patrón de las
corridas base --- la distancia entre los dos modelos depende de la tarea y se
estrecha casi hasta desaparecer en clasificación de temas --- solo que ahora
BERT también gana en AG News, aunque por un margen que sigue dentro del ruido.

\paragraph{Un hallazgo secundario que merece atención.} La cabeza lineal
\emph{mejora} a BERT respecto de su propia corrida base: \num{0,9289} frente a
\num{0,9232} de \textit{accuracy} en SST-2. La corrida base usa la cabeza por
defecto de HuggingFace, que en BERT pasa el \texttt{[CLS]} por un
\textit{pooler} con \texttt{tanh} más \textit{dropout}. Es decir: ese
preprocesado adicional no aporta nada en estas tareas y en SST-2 incluso
estorba. Es coherente con la conclusión del estudio de ablación --- la
capacidad que importa está en el \textit{encoder} --- y llevado un paso más
allá: no solo las capas extra de la cabeza son innecesarias, sino que la
transformación no lineal por defecto puede ser contraproducente.

\section{Limitaciones}
\label{sec:limitaciones}

Cuatro, en orden de importancia para la interpretación de los resultados.

\paragraph{Una sola semilla.} Todas las corridas usan la semilla 42. Sin
repeticiones no hay estimación de la varianza, y por eso a lo largo del
informe las diferencias por debajo de unas décimas se describen como empates
en lugar de como victorias. La afirmación ``la cabeza no importa'' se apoya en
que cinco configuraciones caben en \num{0,0014} de F1, lo cual es consistente
con el ruido esperable pero no es una prueba formal. Repetir con tres o cinco
semillas y reportar media y desviación sería la mejora más valiosa a este
trabajo, y es barata: unas \num{7,5} horas de GPU.

\paragraph{Presupuesto de dos épocas.} El artículo de BERT recomienda entre 2
y 4 épocas para \textit{fine-tuning}. Con 2 se favorece ligeramente al modelo
que converge más rápido, y las curvas de SST-2 sugieren que ahí BERT aún
tenía margen. Es un presupuesto idéntico para los dos modelos, así que no
invalida la comparación, pero sí acota su alcance.

\paragraph{Yelp submuestreado.} Se usan \num{100000} de \num{560000} ejemplos
de train. Los valores absolutos de Yelp están por debajo de lo alcanzable con
el conjunto completo; las diferencias \emph{entre} modelos, que es lo que
interesa, no deberían verse afectadas, porque ambos reciben exactamente el
mismo subconjunto.

\paragraph{Registro de validación incompleto en las corridas base.} Cinco de
las seis corridas base se ejecutaron con una revisión anterior del código de
entrenamiento que no serializaba el resumen de validación por época
(\texttt{desempeno\_val}) en el JSON. La \emph{selección} del mejor
\textit{checkpoint} por F1 de validación sí se aplicó en todas --- está en
las tres revisiones del código ---, y la validación se evaluó durante el
entrenamiento, como demuestran los puntos de \texttt{val\_loss} registrados en
el historial. Lo que falta es solo el registro agregado, motivo por el cual la
Tabla~\ref{tab:desempeno} no incluye columna de F1 de validación y la
Tabla~\ref{tab:mejor}, que proviene de corridas posteriores, sí.
